% RESUMO

\begin{resumo}[RESUMO (OBRIGATÓRIO)]
\begin{Spacing}{1.5}

Uma competência fundamental na formação em Computação é a capacidade de compreender, implementar e selecionar estruturas de dados adequadas para diferentes problemas. Nesse contexto, o avanço de ferramentas de inteligência artificial generativa, como o ChatGPT, tem provocado discussões sobre seus possíveis efeitos na aprendizagem de estudantes em disciplinas de programação. Este trabalho investiga o impacto do uso do ChatGPT na aprendizagem de alunos da disciplina de Algoritmos e Estruturas de Dados, com foco nos conteúdos de Heap e Tabelas Hash. A pesquisa toma como base metodológica um estudo experimental aplicado a estruturas de dados em contexto de graduação, adaptando seu recorte para os tópicos selecionados neste trabalho. A proposta envolve a realização de atividades práticas com grupos de estudantes submetidos a diferentes condições de uso do ChatGPT, permitindo comparar evidências de aprendizagem por meio de questionários, desempenho acadêmico e análise de similaridade de código. Com isso, busca-se observar se o uso da ferramenta influencia a compreensão conceitual e prática dos conteúdos, bem como se contribui para maior semelhança entre as soluções produzidas. Espera-se que os resultados contribuam para a discussão sobre o uso pedagógico de inteligência artificial generativa em cursos de Computação e para o planejamento de atividades avaliativas mais adequadas ao cenário atual.
\end{Spacing}
\vspace{1cm}
\textbf{Palavras-chave}: ChatGPT, inteligência artificial generativa, algoritmos e estruturas de dados, Heap, tabelas hash, aprendizagem.

\end{resumo}

% OBSERVAÇÕES:

% Escolha de 3 a 5 palavras ou termos que descrevam bem o seu trabalho .
% As palavras-chave são separadas por vírgula.

